% ============================================================
% Chapter 6: Experimental Design
%   Section: Selection and Justification of the Evaluation Dataset
%     Subsection: Domain Limitation and a Custom Validation Set
% Included from chapters/06-experimental-design/02-evaluation-dataset/evaluation-dataset.tex
% ============================================================

\subsection{Domain Limitation and a Custom Validation Set}
\label{subsec:custom-gold-set}

Nevertheless, X-Fact does not fully capture the nature of this
work's application domain either. Its claims originate predominantly
from political discourse and controversial public statements, while
the developed system is applied to editorial articles from a
positive-news outlet, with claims combining verifiable facts
(figures, dates, events) with analytical or cultural
interpretations --- a type of claim virtually absent from X-Fact and
which, as documented in Chapter~\ref{ch:results}, exposes a real
limitation of the system in its current version.

For this reason, the quantitative evaluation in this work relies on
\textbf{two complementary sources}: X-Fact \cite{gupta2021xfact}, as
a standardised benchmark enabling comparison of the proposed method
against the state of the art in multilingual, open-retrieval fact
verification; and a custom validation set, built from claims
extracted from real articles published by the outlet under study and
manually labelled, specifically designed to expose and measure the
system's behaviour on interpretive claims --- the real use case for
which the system was conceived, and one that no existing academic
benchmark adequately covers.
